\section{Complete Measured Results}
\label{app:results}
This appendix preserves all operating points from the completed COCO frontier,
the 192-example MileBench E6 lower-memory campaign, the earlier 256-example
MileBench fixed-budget campaign, and the 1,218-example PrefixKV campaign.
R means local reimplementation; MosaicKV exact and Prompt Top-K are locally
measured methods. None of these tables contains external paper-reported
numbers. FullKV is a quality reference and its latency is omitted.
ROUGE-L is scaled by 100, as explicitly labeled; raw metric exports retain
the native $[0,1]$ values.

\subsection{COCO: every shared-byte point}
\begin{table*}[t]
\caption{COCO Caption, 256 images. Each method uses the same measured
initial cap and final/peak storage per sample. CIDEr is a corpus score, not
an average of shard scores. All method--sample rows completed.}
\centering\small% Generated from hash-verified measured artifacts by scripts/export_results.py.
\begin{tabular}{rlrrr}
\toprule
Cap & Method & CIDEr & Final MiB & E2E s\\\midrule
10 & Prompt Top-K & 103.24 & 17.875 & 14.99\\
10 & VL-Cache (R) & 91.83 & 17.875 & 17.75\\
10 & MosaicKV exact & 103.16 & 17.875 & 14.33\\
20 & Prompt Top-K & 108.13 & 27.875 & 15.84\\
20 & VL-Cache (R) & 103.62 & 27.875 & 18.70\\
20 & MosaicKV exact & 107.97 & 27.875 & 15.03\\
40 & Prompt Top-K & 111.59 & 47.875 & 17.22\\
40 & VL-Cache (R) & 107.84 & 47.875 & 20.75\\
40 & MosaicKV exact & 111.48 & 47.875 & 16.68\\
80 & Prompt Top-K & 110.72 & 87.875 & 20.06\\
80 & VL-Cache (R) & 110.04 & 87.875 & 24.28\\
80 & MosaicKV exact & 110.79 & 87.875 & 19.52\\
\bottomrule
\end{tabular}

\end{table*}
The common FullKV reference scores 111.2424 CIDEr, with 285.6064 MiB
mean initial and 293.4814 MiB final live cache. The 64-token output schedule
produces 7.875 MiB decode growth. At 10, 20, 40, and 80 MiB initial,
MosaicKV--VL-Cache CIDEr differences are respectively 11.34, 4.35, 3.64,
and 0.75; paired 95\% intervals are $[6.55,16.49]$, $[0.01,8.90]$,
$[0.53,7.06]$, and $[-3.09,4.33]$. The last point is unresolved.
At 40 MiB, Prompt Top-K's 111.59 exceeds MosaicKV's 111.48 in point
estimate, with an unresolved paired difference. Results therefore distinguish
the published-method port from the strongest simple selector.

\subsection{MileBench E6: every fixed-budget point}
\begin{table*}[t]
\caption{MileBench E6: 192 unique examples, 48 per declared subtask; 96
applicable examples per metric family. Final storage equals peak storage.
All 3,072 rows completed with no fallback, target shortfall, or duplicate.}
\centering\small% Generated from hash-verified measured artifacts by scripts/export_results.py.
\begin{tabular}{rlrrrr}
\toprule
Cap & Method & Acc. (\%) & R-L ($\times100$) & Final MiB & E2E s\\\midrule
16 & Prompt Top-K & 20.83 & 14.47 & 47.500 & 16.10\\
16 & LOOK-M (R) & 3.12 & 12.31 & 47.500 & 24.36\\
16 & MosaicKV exact & 21.88 & 14.52 & 47.500 & 14.95\\
20 & Prompt Top-K & 20.83 & 14.92 & 51.500 & 16.39\\
20 & LOOK-M (R) & 12.50 & 14.25 & 51.500 & 26.35\\
20 & MosaicKV exact & 21.88 & 14.78 & 51.500 & 15.34\\
28 & Prompt Top-K & 21.88 & 15.67 & 59.500 & 17.13\\
28 & LOOK-M (R) & 21.88 & 14.11 & 59.500 & 30.04\\
28 & MosaicKV exact & 20.83 & 15.57 & 59.500 & 15.84\\
40 & Prompt Top-K & 22.92 & 15.88 & 71.500 & 18.03\\
40 & LOOK-M (R) & 25.00 & 14.48 & 71.500 & 34.22\\
40 & MosaicKV exact & 22.92 & 15.85 & 71.500 & 16.71\\
80 & Prompt Top-K & 23.96 & 15.96 & 111.500 & 20.91\\
80 & LOOK-M (R) & 21.88 & 15.67 & 111.500 & 47.42\\
80 & MosaicKV exact & 23.96 & 16.17 & 111.500 & 19.74\\
\bottomrule
\end{tabular}

\end{table*}
\begin{table*}[t]
\caption{E6 subtask scores at 40 MiB; 48 examples each. Temporal rows use
Accuracy (\%); semantic rows use ROUGE-L $\times100$. The columns are
not averaged across metric families.}
\centering\small% Generated from hash-verified measured artifacts by scripts/export_results.py.
\begin{tabular}{lrrrr}
\toprule
Subtask & FullKV & Top-K & LOOK-M (R) & MosaicKV\\\midrule
ActionPrediction & 20.83 & 20.83 & 20.83 & 20.83\\
ActionSequence & 25.00 & 25.00 & 29.17 & 25.00\\
CLEVR{-}Change & 14.67 & 15.81 & 14.19 & 15.52\\
Spot{-}the{-}Diff & 17.26 & 15.96 & 14.77 & 16.17\\
\bottomrule
\end{tabular}

\end{table*}
The FullKV category references are 22.9167\% Accuracy and 15.9631 ROUGE-L
points. Its mean initial storage is 315.4688 MiB and final storage is
346.9688 MiB. Accuracy is low even without compression; this limits the
precision and breadth of conclusions about high-level temporal reasoning.

\subsection{MileBench: separate earlier 256-example cohort}
\begin{table*}[t]
\caption{Earlier fixed-budget cohort: 64 examples per subtask, 256 total,
128 per metric family. Its prior use is disclosed. These measurements are
not pooled with E6 or relabeled as the same cohort.}
\centering\small% Generated from hash-verified measured artifacts by scripts/export_results.py.
\begin{tabular}{rlrrrr}
\toprule
Cap & Method & Acc. (\%) & R-L ($\times100$) & Final MiB & E2E s\\\midrule
10 & Prompt Top-K & 25.00 & 15.16 & 41.500 & 15.65\\
10 & LOOK-M (R) & 0.78 & 9.81 & 41.500 & 21.49\\
10 & MosaicKV exact & 26.56 & 15.31 & 41.500 & 14.61\\
20 & Prompt Top-K & 24.22 & 15.87 & 51.500 & 16.36\\
20 & LOOK-M (R) & 13.28 & 14.95 & 51.500 & 26.32\\
20 & MosaicKV exact & 24.22 & 16.02 & 51.500 & 15.34\\
40 & Prompt Top-K & 30.47 & 16.97 & 71.500 & 18.01\\
40 & LOOK-M (R) & 30.47 & 15.83 & 71.500 & 34.12\\
40 & MosaicKV exact & 30.47 & 16.82 & 71.500 & 16.76\\
80 & Prompt Top-K & 31.25 & 17.22 & 111.500 & 20.86\\
80 & LOOK-M (R) & 33.59 & 16.14 & 111.500 & 48.00\\
80 & MosaicKV exact & 30.47 & 17.02 & 111.500 & 19.95\\
\bottomrule
\end{tabular}

\end{table*}
At 40 MiB, the earlier cohort's MosaicKV--LOOK-M ROUGE-L gain is 0.9846
points, interval $[0.0075,2.0091]$; Accuracy ties in point estimate with
interval $[-5.4688,5.4688]$ percentage points. At 80 MiB LOOK-M's Accuracy
point estimate exceeds MosaicKV by 3.125 points. The earlier FullKV category
scores are 32.8125\% Accuracy and 16.8532 ROUGE-L points, demonstrating
why cross-cohort point estimates cannot be treated as a single curve.

\subsection{Registered lower-memory tests}
\begin{table*}[t]
\caption{All feasible E6 lower-memory comparisons. B/M indicates baseline
initial cap / MosaicKV initial cap. Each interval is the paired quality
difference, MosaicKV minus baseline, using 98.75\% per-comparison confidence
for the registered family-wise correction. Quality margin is zero.
Memory passes at initial, final, and peak boundaries; quality does not.}
\centering\scriptsize\resizebox{\textwidth}{!}{% Generated from hash-verified measured artifacts by scripts/export_results.py.
\begin{tabular}{lllrrrl}
\toprule
Baseline & B/M MiB & Metric & Init. reduction \% & Final/peak \% & Adjusted CI ($\times100$) & Gate\\\midrule
Prompt Top-K & 40/16 & ROUGE-L & 60.0 & 33.57 & -2.60, -0.14 & Fail\\
Prompt Top-K & 40/16 & Accuracy & 60.0 & 33.57 & -4.17, 0.00 & Fail\\
Prompt Top-K & 80/40 & ROUGE-L & 50.0 & 35.87 & -1.36, 0.99 & Fail\\
Prompt Top-K & 80/40 & Accuracy & 50.0 & 35.87 & -5.21, 3.12 & Fail\\
LOOK-M (R) & 40/16 & ROUGE-L & 60.0 & 33.57 & -1.61, 1.71 & Fail\\
LOOK-M (R) & 40/16 & Accuracy & 60.0 & 33.57 & -9.38, 2.08 & Fail\\
LOOK-M (R) & 80/40 & ROUGE-L & 50.0 & 35.87 & -1.03, 1.38 & Fail\\
LOOK-M (R) & 80/40 & Accuracy & 50.0 & 35.87 & -4.17, 7.29 & Fail\\
\bottomrule
\end{tabular}
}
\end{table*}
Baseline caps 16, 20, and 28 MiB have no eligible smaller MosaicKV point on
the registered grid after the 31.5 MiB decode growth is included. They are
recorded as physically infeasible comparisons, not omitted or interpolated.
Both metric families must pass on the same pair for a lane-wide claim.
The 16-example execution validation produced some metric-specific positives;
the larger 192-example campaign did not confirm the joint claim. The larger
result takes precedence; execution success is distinct from quality success.

The COCO 20-versus-40 MiB example is descriptive only: candidate CIDEr
107.9651 versus baseline 107.8443 and final storage 27.875 versus 47.875 MiB.
Its 41.78\% reduction is measured; equal-quality inference is not established
by those point estimates. A zero-margin confidence test cannot be replaced
by observing a positive mean. Prompt Top-K's 20 MiB score is also 108.1280.

\subsection{PrefixKV: every native and common-byte result}
\begin{table*}[t]
\caption{PrefixKV campaign: 1,000 descriptions and 218 paper-style MM-Vet
questions, 8,526 rows, no failed or duplicate rows. Native is BF16 at nominal
ratio 0.5 with method-specific maintenance. The separate 40 MiB lane is FP16
with append-only compressed decoding. Its FullKV reference is distinct.}
\centering\scriptsize\resizebox{\textwidth}{!}{% Generated from hash-verified measured artifacts by scripts/export_results.py.
\begin{tabular}{lllrrrrr}
\toprule
Task & Lane & Method & PPL $\downarrow$ & R-L $\uparrow$ & Initial MiB & Peak MiB & E2E s\\\midrule
LLaVA-Desc. & Native & FullKV & 2.9555 & 100.00 & 298.83 & 810.33 & ---\\
LLaVA-Desc. & Native & PrefixKV (R) & 3.0889 & 63.42 & 149.70 & 406.42 & 52.96\\
LLaVA-Desc. & Native & MosaicKV exact & 2.9620 & 71.09 & 102.62 & 614.12 & 63.91\\
MM-Vet & Native & FullKV & 5.2778 & 100.00 & 301.42 & 812.92 & ---\\
MM-Vet & Native & PrefixKV (R) & 5.3339 & 64.97 & 151.01 & 407.74 & 52.96\\
MM-Vet & Native & MosaicKV exact & 5.2775 & 79.03 & 103.63 & 615.13 & 64.03\\
LLaVA-Desc. & 40 MiB & FullKV & 2.9540 & 100.00 & 298.83 & 810.33 & ---\\
LLaVA-Desc. & 40 MiB & Prompt Top-K & 3.0017 & 63.20 & 40.00 & 551.50 & 59.37\\
LLaVA-Desc. & 40 MiB & PrefixKV (R) & 3.3780 & 13.99 & 40.00 & 551.50 & 53.23\\
LLaVA-Desc. & 40 MiB & MosaicKV exact & 3.0033 & 63.09 & 40.00 & 551.50 & 58.93\\
MM-Vet & 40 MiB & FullKV & 5.2714 & 100.00 & 301.42 & 812.92 & ---\\
MM-Vet & 40 MiB & Prompt Top-K & 5.2735 & 72.62 & 40.00 & 551.50 & 59.46\\
MM-Vet & 40 MiB & PrefixKV (R) & 5.7577 & 13.52 & 40.00 & 551.50 & 53.23\\
MM-Vet & 40 MiB & MosaicKV exact & 5.2726 & 72.71 & 40.00 & 551.50 & 58.93\\
\bottomrule
\end{tabular}
}
\end{table*}
All 153 GPU shard jobs completed. The original analyzer wrote numerical
Parquet and LaTeX outputs before failing on a missing optional plotting
package. Paper export rechecks 8,526 rows, exactly 1,000/218 per arm,
zero duplicates/failures, and equal per-sample initial/final/peak bytes in
the shared lane. This is independent numerical verification, not a fabricated
successful original report.

The native PrefixKV path uses the pinned 0.5 profile, upstream rounding,
protected boundaries, and profile-ratio decode maintenance. Four-example
official/local parity evidence is preserved separately. It does not establish
a universal checkpoint/environment equivalence or a numerical reproduction of
an unregistered published ROUGE-L target. Common-byte rescaling and append-only
maintenance are explicit adaptations; their output collapse is evidence about
that adaptation regime, not an intrinsic failure of original PrefixKV.

\section{Interpreting Time and Cache Growth}
End-to-end observations include planning. On COCO at 40 MiB, only 0.1724 s
of the 4.0776 s MosaicKV--VL-Cache reduction lies in decode. At equal bytes,
cache volume alone cannot explain that difference. A missing upstream fused
kernel, differing packing paths, and host-side planning affect local ports.
No result here implies superiority over authors' optimized kernel timings.

At 1,024 generated tokens, the common PrefixKV lane grows by 511.5 MiB;
a 40 MiB initial saving is therefore a small share of the final cache.
Maintaining the native eviction policy changes both bytes and output behavior.
The manuscript reports both, rather than describing a nominal retention
comparison as equal physical capacity.

\section{Exploratory Algorithms and Inactive Extensions}
Prompt-output reconstruction, bounded corrective swaps, draft-query vetoes,
joint draft ranking, and short-horizon logit checks were explored during
development. They address a real limitation of prompt attention: reconstruction
error can improve without preserving autoregressive answers. Attention matching
is related prior work \citep{attentionmatching2026}. These development
candidates are not the frozen 1/256-routed policy used in the measured tables;
their inspected cohorts cannot serve as untouched validation of that policy.
No exploratory algorithm receives a measured gain by association with the
completed campaigns. Likewise, prototype and residual-repair specifications
in the following sections are architectural extensions rather than evaluated
performance claims.
